\begin{table}[t]
\caption{Research questions and evidence. The supplement retains the full experimental record.}\label{tab:questions}
\centering\small\setlength{\tabcolsep}{3pt}\begin{tabular}{p{.25\textwidth}p{.32\textwidth}p{.34\textwidth}}\toprule
Question & Evidence & Claim assessed\\\midrule
Behavioral model choice & Frozen pre-crisis response experiment & Static, predictive-selected and retained-load models choose different actions.\\
Source of the difference & Saved policy audit, simple selector and true-generator planning & Separate action choice, threshold choice and model error.\\
Operating assumptions & Effect, event-definition and horizon sensitivity; discovery check & Identify the conditions covered by the response evaluation.\\
Observed workload signal & Public BPI prediction checks & Workload may aid prediction; this does not evaluate real interventions.\\\bottomrule
\end{tabular}\end{table}

\subsection{Diagnostic analysis and supporting evidence}
After completing the primary experiment, we fixed a separate diagnostic protocol. It adds a simple selector that compares estimated service and rework amplification between the training-data 10th and 90th percentiles of load, then chooses the corresponding action at threshold 20. The scores are log ratios of expected B service duration and of the expected repetition multiplier $1/(1-p)$, with negative scores set to zero; complexity is held at zero. It chooses rework when the rework score is larger and capacity otherwise. This rule uses fitted parameters and training states only.

A true-generator planner uses the original planning scenarios and repetitions with the generating equations. Its choice is saved before scoring the existing independent test actions. It is an oracle for model knowledge within the finite policy catalog, not a globally optimal or clairvoyant controller. We also compute the best fixed policy on the realized test data as a separately labelled hindsight lower envelope. Their difference from the deployed candidate measures distinct sources of decision loss. Later diagnostic intervals are descriptive, not additional confirmatory endpoints.

The intervention-effect check keeps the original decisions fixed under assumed rework reduction 50\% and evaluates actual reductions 30\% and 70\%. Other checks reuse stored trajectories at WIP thresholds 80/100/120, sustained durations 100/200/300, horizons 1400/1600/1800 and terminal unfinished-case weights 0/100/200. These alter evaluation assumptions without selecting policies on test outcomes. A discovery integration check uses only training traces and explicitly separates the recovered serial family from earlier parallel-discovery errors.

Earlier simulation campaigns assess trajectory prediction and structural fidelity. Public BPI 2012 and BPI 2017 logs~\cite{bpi12,bpi17} provide chronological tests of workload features for duration and next-start prediction. These checks establish observational signal, rather than validate intervention effects; their complete methods, controls and results are in the supplement.
